% Generated draft: requires mathematical and source review.
\documentclass[11pt]{article}
\usepackage[margin=25mm]{geometry}
\usepackage{fontspec}
\setmainfont{Latin Modern Roman}
\usepackage{amsmath,amssymb,mathtools,unicode-math}
\setmathfont{Latin Modern Math}
\usepackage{enumitem,needspace,xurl,hyperref}
\setlength{\parindent}{0pt}
\setlength{\parskip}{5pt}
\setlength{\emergencystretch}{3em}
\newcommand{\N}{\mathbb{N}}
\newcommand{\Z}{\mathbb{Z}}
\newcommand{\Q}{\mathbb{Q}}
\newcommand{\R}{\mathbb{R}}
\newcommand{\C}{\mathbb{C}}
\newcommand{\F}{\mathbb{F}}
\newcommand{\A}{\mathbb{A}}
\newcommand{\PP}{\mathbb{P}}
\newcommand{\E}{\mathbb{E}}
\newcommand{\eps}{\varepsilon}
\newcommand{\dd}{\,\mathrm{d}}
\begin{document}
\section*{UnsolvedMath ID: 30005366}
\textbf{Source problem code:} catalogue entry
\subsection{Definitions and mathematical statement}
Let \(M\) be a closed, connected, oriented, irreducible rational homology \(3\)-sphere. Here rational homology sphere means that \(H_k(M;\Q)\cong H_k(S^3;\Q)\) for every \(k\), and irreducibility means that every embedded \(2\)-sphere bounds a \(3\)-ball. Write \(S^1=\R/\Z\), and let \(\operatorname{Homeo}^{+}(S^1)\) and \(\operatorname{Homeo}^{+}(\R)\) denote the groups of orientation-preserving homeomorphisms. Define
\[
\operatorname{Homeo}_{\Z}(\R)=\{f\in\operatorname{Homeo}^{+}(\R): f(x+1)=f(x)+1\text{ for every }x\in\R\}.
\]
The supplied exposition identifies this group as the universal covering group of \(\operatorname{Homeo}^{+}(S^1)\). Its projection \(p\) is given by \(p(f)(x+\Z)=f(x)+\Z\).

For a nontrivial group homomorphism
\[
\rho:\pi_1(M)\longrightarrow\operatorname{Homeo}^{+}(S^1),
\]
a lift means a group homomorphism \(\widetilde\rho:\pi_1(M)\to\operatorname{Homeo}_{\Z}(\R)\) satisfying \(p\circ\widetilde\rho=\rho\). The source describes an Euler class \(e(\rho)\in H^2(M)\) as the obstruction to such a lift. It does not specify the cohomology coefficients; under the usual integral Euler-class convention, this notation would mean \(e(\rho)\in H^2(M;\Z)\). This coefficient interpretation requires confirmation.

A group \(G\) is left-orderable if it admits a strict total order \(<\) such that \(g<h\) implies \(ag<ah\) for every \(a,g,h\in G\). The excerpt discusses representations arising from foliations and flows on \(M\), their use in establishing left-orderability of \(\pi_1(M)\) in unspecified situations, and analogous existence results for taut foliations in connection with the L-space conjecture.

\textbf{Unresolved mathematical target.} The supplied record contains introductory exposition, with no standalone question, conjecture, or specified implication to formalize. It gives neither the additional hypotheses for the situations mentioned nor an exact formulation involving taut foliations or the L-space conjecture. Accordingly, the definitions and background above do not constitute a recovered open problem. The supplied verification field labels the statement unrecoverable and says that specific conjectures occur in adjacent records, whose contents are not supplied here.
\subsection{Short English statement}
For a closed, connected, oriented, irreducible rational homology \(3\)-sphere \(M\), the excerpt introduces circle actions of \(\pi_1(M)\) and the Euler-class obstruction to lifting them to \(\operatorname{Homeo}_{\Z}(\R)\). It discusses connections with left-orderability and taut foliations but states no precise problem. The supplied open label reflects an unclear-claim curation rule, not a verified mathematical open-status assessment.
\subsection{Sources}
\begin{itemize}
\item Local UnsolvedMath registry, permanent ID \(30005366\), title \emph{Background on Circle Actions and the L-Space Conjecture}; Oberwolfach Reports Open Problems collection, exact source code \texttt{catalogue entry}. The supplied statement status is \texttt{unrecoverable}; no clean statement is supplied.
\item \emph{Low-dimensional Topology} (2023), Oberwolfach Reports, pp.\ 151--213, DOI \texttt{10.4171/owr/2023/3}: \url{https://doi.org/10.4171/owr/2023/3}. This is the supplied report citation and the literature-source link labeled ``original Oberwolfach report.''
\item The excerpt attributes the contribution \emph{\(\operatorname{Homeo}^{+}(S^1)\) representations and the L-space conjecture} to Steven Boyer, joint work with Cameron McA. Gordon and Ying Hu. No separate problem proposer is supplied.
\item Literature assessment dated exactly \texttt{2026-08-22}: supplied dated curation, not independently verified here. It classifies the record as open under an unclear-claim rule because the excerpt is background rather than a complete question. The background requests independent verification of time-sensitive status and cited claims.
\item The supplied statement-verification field reports that the official report confirms the introductory nature of the extraction and that adjacent records contain the specific conjectures. This is a supplied verification claim, not an independent source consultation for this draft.
\end{itemize}
\end{document}
